\section{Commutative rings and modules}

\subsection{Reminder: Commutative rings}

\begin{definition}
\label{def.alg.commring}
\lean{AlgebraicCombinatorics.FPS.commRing_add_comm, AlgebraicCombinatorics.FPS.commRing_add_assoc, AlgebraicCombinatorics.FPS.commRing_add_zero, AlgebraicCombinatorics.FPS.commRing_sub_iff_add, AlgebraicCombinatorics.FPS.commRing_mul_comm, AlgebraicCombinatorics.FPS.commRing_mul_assoc, AlgebraicCombinatorics.FPS.commRing_left_distrib, AlgebraicCombinatorics.FPS.commRing_mul_one, AlgebraicCombinatorics.FPS.commRing_mul_zero}
\leantarget
\leanok
A \emph{commutative ring} means a set $K$ equipped
with three maps%
\begin{align*}
\oplus &  :K\times K\rightarrow K,\\
\ominus &  :K\times K\rightarrow K,\\
\odot &  :K\times K\rightarrow K
\end{align*}
and two elements $\mathbf{0}\in K$ and $\mathbf{1}\in K$ satisfying the
following axioms:

\begin{enumerate}
\item \emph{Commutativity of addition:} We have $a\oplus b=b\oplus a$ for all
$a,b\in K$.

\item \emph{Associativity of addition:} We have $a\oplus\left(  b\oplus
c\right)  =\left(  a\oplus b\right)  \oplus c$ for all $a,b,c\in K$.

\item \emph{Neutrality of zero:} We have $a\oplus\mathbf{0}=\mathbf{0}\oplus
a=a$ for all $a\in K$.

\item \emph{Subtraction undoes addition:} Let $a,b,c\in K$. We have $a\oplus
b=c$ if and only if $a=c\ominus b$.

\item \emph{Commutativity of multiplication:} We have $a\odot b=b\odot a$ for
all $a,b\in K$.

\item \emph{Associativity of multiplication:} We have $a\odot\left(  b\odot
c\right)  =\left(  a\odot b\right)  \odot c$ for all $a,b,c\in K$.

\item \emph{Distributivity:} We have%
\[
a\odot\left(  b\oplus c\right)  =\left(  a\odot b\right)  \oplus\left(  a\odot
c\right)  \ \ \ \ \ \ \ \ \ \ \text{and}\ \ \ \ \ \ \ \ \ \ \left(  a\oplus
b\right)  \odot c=\left(  a\odot c\right)  \oplus\left(  b\odot c\right)
\]
for all $a,b,c\in K$.

\item \emph{Neutrality of one:} We have $a\odot\mathbf{1}=\mathbf{1}\odot a=a$
for all $a\in K$.

\item \emph{Annihilation:} We have $a\odot\mathbf{0}=\mathbf{0}\odot
a=\mathbf{0}$ for all $a\in K$.
\end{enumerate}

The operations $\oplus$, $\ominus$ and $\odot$ are called the \emph{addition},
the \emph{subtraction} and the \emph{multiplication} of the ring $K$.
When confusion is unlikely,
we will denote these three operations $\oplus$, $\ominus$ and $\odot$ by $+$,
$-$ and $\cdot$, respectively, and we will abbreviate $a\odot b=a\cdot b$ by
$ab$.

The elements $\mathbf{0}$ and $\mathbf{1}$ are called the \emph{zero} and the
\emph{unity} (or the \emph{one}) of the ring $K$.
We will simply call these elements $0$ and $1$ when confusion with the
corresponding numbers is unlikely.

We will use \emph{PEMDAS conventions} for the three operations $\oplus$,
$\ominus$ and $\odot$. These imply that the operation $\odot$ has higher
precedence than $\oplus$ and $\ominus$, while the operations $\oplus$ and
$\ominus$ are left-associative.
\end{definition}

\subsection{Standard rules in commutative rings}

\begin{lemma}
\label{lem.commring.neg-add}
\lean{AlgebraicCombinatorics.FPS.neg_add_distrib}
\leanhelper
\leanok
Let $K$ be a commutative ring. For any $a,b\in K$, we have
$-\left(  a+b\right)  =\left(  -a\right)  +\left(  -b\right)$.
\end{lemma}

\begin{proof}
\leanok
Standard ring theory.
\end{proof}

\begin{lemma}
\label{lem.commring.neg-neg}
\lean{AlgebraicCombinatorics.FPS.neg_neg_eq}
\leanhelper
\leanok
Let $K$ be a commutative ring. For any $a\in K$, we have
$-\left(  -a\right)  =a$.
\end{lemma}

\begin{proof}
\leanok
Standard ring theory.
\end{proof}

\begin{lemma}
\label{lem.commring.add-zsmul}
\lean{AlgebraicCombinatorics.FPS.add_zsmul_distrib}
\leanhelper
\leanok
Let $K$ be a commutative ring. For any $a\in K$ and $n,m\in\mathbb{Z}$, we have
$\left(  n+m\right)  a=na+ma$.
\end{lemma}

\begin{proof}
\leanok
Standard ring theory.
\end{proof}

\begin{lemma}
\label{lem.commring.mul-zsmul}
\lean{AlgebraicCombinatorics.FPS.mul_zsmul_assoc}
\leanhelper
\leanok
Let $K$ be a commutative ring. For any $a\in K$ and $n,m\in\mathbb{Z}$, we have
$\left(  nm\right)  a=n\left(  ma\right)$.
\end{lemma}

\begin{proof}
\leanok
Standard ring theory.
\end{proof}

\begin{lemma}
\label{lem.commring.mul-sub}
\lean{AlgebraicCombinatorics.FPS.mul_sub_distrib}
\leanhelper
\leanok
Let $K$ be a commutative ring. For any $a,b,c\in K$, we have
$a\left(  b-c\right)  =\left(  ab\right)  -\left(  ac\right)$.
\end{lemma}

\begin{proof}
\leanok
Standard ring theory.
\end{proof}

\begin{lemma}
\label{lem.commring.mul-pow}
\lean{AlgebraicCombinatorics.FPS.mul_pow_eq}
\leanhelper
\leanok
Let $K$ be a commutative ring. For any $a,b\in K$ and $n\in\mathbb{N}$, we have
$\left(  ab\right)  ^{n}=a^{n}b^{n}$.
\end{lemma}

\begin{proof}
\leanok
Standard ring theory (induction on $n$).
\end{proof}

\begin{lemma}
\label{lem.commring.pow-add}
\lean{AlgebraicCombinatorics.FPS.pow_add_eq}
\leanhelper
\leanok
Let $K$ be a commutative ring. For any $a\in K$ and $n,m\in\mathbb{N}$, we have
$a^{n+m}=a^{n}a^{m}$.
\end{lemma}

\begin{proof}
\leanok
Standard ring theory (induction on $m$).
\end{proof}

\begin{lemma}
\label{lem.commring.pow-mul}
\lean{AlgebraicCombinatorics.FPS.pow_mul_eq}
\leanhelper
\leanok
Let $K$ be a commutative ring. For any $a\in K$ and $n,m\in\mathbb{N}$, we have
$a^{nm}=\left(  a^{n}\right)  ^{m}$.
\end{lemma}

\begin{proof}
\leanok
Standard ring theory (induction on $m$).
\end{proof}

\begin{theorem}
\label{thm.commring.binomial}
\lean{AlgebraicCombinatorics.FPS.binomial_theorem}
\leanhelper
\leanok
\textbf{(Binomial Theorem.)}
Let $K$ be a commutative ring. For any $a,b\in K$ and $n\in\mathbb{N}$,
\[
(a+b)^n = \sum_{k=0}^{n} \binom{n}{k} a^k b^{n-k}.
\]
\end{theorem}

\begin{proof}
Standard; this is the binomial theorem from Mathlib.
\end{proof}

\begin{lemma}
\label{lem.commring.binomial-sub}
\lean{AlgebraicCombinatorics.FPS.binomial_theorem_sub}
\leanhelper
\leanok
Let $K$ be a commutative ring. For any $a,b\in K$ and $n\in\mathbb{N}$,
\[
(a-b)^n = \sum_{m=0}^{n} (-1)^{m+n} \cdot a^m \cdot b^{n-m} \cdot \binom{n}{m}.
\]
\end{lemma}

\begin{proof}
\leanok
Write $a - b = a + (-b)$ and apply Theorem~\ref{thm.commring.binomial}.
\end{proof}


\subsection{Finite sums and products}

\begin{lemma}
\label{lem.commring.sum-disjoint}
\lean{AlgebraicCombinatorics.FPS.sum_disjoint_union}
\leanhelper
\leanok
Let $K$ be a commutative ring. Let $S$ be a finite type, and let
$X, Y$ be disjoint finite subsets of $S$. For any family $(a_s)_{s \in S}$
of elements of $K$, we have
$\sum_{s\in X\cup Y}a_{s}=\sum_{s\in X}a_{s}+\sum_{s\in Y}a_{s}$.
\end{lemma}

\begin{proof}
\leanok
Standard.
\end{proof}

\begin{lemma}
\label{lem.commring.sum-add}
\lean{AlgebraicCombinatorics.FPS.sum_add_distrib'}
\leanhelper
\leanok
Let $K$ be a commutative ring. Let $T$ be a finite set and $(a_s)_{s\in T}$, $(b_s)_{s\in T}$
be two families of elements of $K$. Then
$\sum_{s\in T}\left(  a_{s}+b_{s}\right)  =\sum_{s\in T}a_{s}+\sum_{s\in T}b_{s}$.
\end{lemma}

\begin{proof}
\leanok
Standard.
\end{proof}

\begin{lemma}
\label{lem.commring.sum-empty}
\lean{AlgebraicCombinatorics.FPS.sum_empty}
\leanhelper
\leanok
Let $K$ be a commutative ring. If $S=\varnothing$, then $\sum_{s\in S}a_{s}=0$.
\end{lemma}

\begin{proof}
\leanok
By definition.
\end{proof}

\begin{lemma}
\label{lem.commring.prod-empty}
\lean{AlgebraicCombinatorics.FPS.prod_empty}
\leanhelper
\leanok
Let $K$ be a commutative ring. If $S=\varnothing$, then $\prod_{s\in S}a_{s}=1$.
\end{lemma}

\begin{proof}
\leanok
By definition.
\end{proof}

\subsection{$K$-modules}

\begin{definition}
\label{def.alg.module}
\lean{AlgebraicCombinatorics.FPS.module_add_comm, AlgebraicCombinatorics.FPS.module_add_assoc, AlgebraicCombinatorics.FPS.module_add_zero, AlgebraicCombinatorics.FPS.module_sub_iff_add, AlgebraicCombinatorics.FPS.module_smul_assoc, AlgebraicCombinatorics.FPS.module_smul_add, AlgebraicCombinatorics.FPS.module_add_smul, AlgebraicCombinatorics.FPS.module_one_smul, AlgebraicCombinatorics.FPS.module_zero_smul, AlgebraicCombinatorics.FPS.module_smul_zero}
\leantarget
\leanok
Let $K$ be a commutative ring.

A $K$\emph{-module} means a set $M$ equipped with three maps
\begin{align*}
\oplus &  :M\times M\rightarrow M,\\
\ominus &  :M\times M\rightarrow M,\\
\rightharpoonup &  :K\times M\rightarrow M
\end{align*}
(notice that the third map has domain $K\times M$, not $M\times M$) and an
element $\overrightarrow{0}\in M$ satisfying the following axioms:

\begin{enumerate}
\item \emph{Commutativity of addition:} We have $a\oplus b=b\oplus a$ for all
$a,b\in M$.

\item \emph{Associativity of addition:} We have $a\oplus\left(  b\oplus
c\right)  =\left(  a\oplus b\right)  \oplus c$ for all $a,b,c\in M$.

\item \emph{Neutrality of zero:} We have $a\oplus\overrightarrow{0}%
=\overrightarrow{0}\oplus a=a$ for all $a\in M$.

\item \emph{Subtraction undoes addition:} Let $a,b,c\in M$. We have $a\oplus
b=c$ if and only if $a=c\ominus b$.

\item \emph{Associativity of scaling:} We have $u\rightharpoonup\left(
v\rightharpoonup a\right)  =\left(  uv\right)  \rightharpoonup a$ for all
$u,v\in K$ and $a\in M$.

\item \emph{Left distributivity:} We have $u\rightharpoonup\left(  a\oplus
b\right)  =\left(  u\rightharpoonup a\right)  \oplus\left(  u\rightharpoonup
b\right)  $ for all $u\in K$ and $a,b\in M$.

\item \emph{Right distributivity:} We have $\left(  u+v\right)
\rightharpoonup a=\left(  u\rightharpoonup a\right)  \oplus\left(
v\rightharpoonup a\right)  $ for all $u,v\in K$ and $a\in M$.

\item \emph{Neutrality of one:} We have $1\rightharpoonup a=a$ for all $a\in
M$.

\item \emph{Left annihilation:} We have $0\rightharpoonup a=\overrightarrow{0}%
$ for all $a\in M$.

\item \emph{Right annihilation:} We have $u\rightharpoonup\overrightarrow{0}%
=\overrightarrow{0}$ for all $u\in K$.
\end{enumerate}

The operations $\oplus$, $\ominus$ and $\rightharpoonup$ are called the
\emph{addition}, the \emph{subtraction} and the \emph{scaling} (or the
$K$\emph{-action}) of the $K$-module $M$. When confusion is unlikely, we will
denote these three operations $\oplus$, $\ominus$ and $\rightharpoonup$ by
$+$, $-$ and $\cdot$, respectively, and we will abbreviate $a\rightharpoonup
b=a\cdot b$ by $ab$.

The element $\overrightarrow{0}$ is called the \emph{zero} (or the \emph{zero
vector}) of the $K$-module $M$. We will usually just call it $0$.

When $M$ is a $K$-module, the elements of $M$ are called \emph{vectors}, while
the elements of $K$ are called \emph{scalars}.

We will use \emph{PEMDAS conventions} for the three operations $\oplus$,
$\ominus$ and $\rightharpoonup$, with the operation $\rightharpoonup$ having
higher precedence than $\oplus$ and $\ominus$.
\end{definition}

\subsection{Additive inverses in modules}

\begin{lemma}
\label{lem.module.neg-one-smul}
\lean{AlgebraicCombinatorics.FPS.module_neg_eq_neg_one_smul}
\leanhelper
\leanok
Let $K$ be a commutative ring and $M$ a $K$-module. For any $a\in M$,
the additive inverse of $a$ can be constructed as $(-1)\cdot a$, i.e.,
$-a = (-1)\cdot a$.
\end{lemma}

\begin{proof}
\leanok
Standard module theory.
\end{proof}

\begin{lemma}
\label{lem.module.sub-eq-add-neg}
\lean{AlgebraicCombinatorics.FPS.module_sub_eq_add_neg}
\leanhelper
\leanok
Let $K$ be a commutative ring and $M$ a $K$-module. For any $a,b\in M$,
we have $a - b = a + (-b)$.
\end{lemma}

\begin{proof}
\leanok
Standard module theory.
\end{proof}

\begin{lemma}
\label{lem.module.sub-eq-add-neg-one-smul}
\lean{AlgebraicCombinatorics.FPS.module_sub_eq_add_neg_one_smul}
\leanhelper
\leanok
Let $K$ be a commutative ring and $M$ a $K$-module. For any $a,b\in M$,
we have $a - b = a + (-1)\cdot b$.
\end{lemma}

\begin{proof}
\leanok
Combine the fact that $-b = (-1)\cdot b$ with $a - b = a + (-b)$.
\end{proof}

\begin{lemma}
\label{lem.module.neg-add}
\lean{AlgebraicCombinatorics.FPS.module_neg_add}
\leanhelper
\leanok
Let $K$ be a commutative ring and $M$ a $K$-module. For any $a,b\in M$,
we have $-(a+b) = (-a) + (-b)$.
\end{lemma}

\begin{proof}
\leanok
Standard module theory.
\end{proof}

\begin{lemma}
\label{lem.module.neg-neg}
\lean{AlgebraicCombinatorics.FPS.module_neg_neg}
\leanhelper
\leanok
Let $K$ be a commutative ring and $M$ a $K$-module. For any $a\in M$,
we have $-(-a) = a$.
\end{lemma}

\begin{proof}
\leanok
Standard module theory.
\end{proof}

\subsection{Inverses in commutative rings}

\begin{definition}
\label{def.commring.inverse}
\lean{AlgebraicCombinatorics.FPS.IsInverse}
\leanhelper
\leanok
Let $L$ be a commutative ring and $a, b \in L$.

\textbf{(a)} We say that $b$ is an \emph{inverse} (or \emph{multiplicative inverse})
of $a$ if $a \cdot b = 1$.

\textbf{(b)} We say that $a$ is \emph{invertible} in $L$ (or a \emph{unit} of $L$)
if $a$ has an inverse.
\end{definition}

\begin{lemma}
\label{lem.commring.inverse-uni}
\lean{AlgebraicCombinatorics.FPS.isInverse_unique}
\leanhelper
\leanok
Let $L$ be a commutative ring. If $b$ and $c$ are both inverses of $a$,
then $b = c$. In other words, the inverse of an element (if it exists) is unique.
\end{lemma}

\begin{proof}
\leanok
We have $b = b \cdot 1 = b \cdot (a \cdot c) = (b \cdot a) \cdot c = 1 \cdot c = c$.
\end{proof}

\begin{definition}
\label{def.commring.invertible}
\lean{AlgebraicCombinatorics.FPS.IsInvertible}
\leanhelper
\leanok
An element $a \in L$ is \emph{invertible} if there exists $b \in L$ such that
$a \cdot b = 1$.
\end{definition}

\begin{lemma}
\label{lem.commring.invertible-iff-isUnit}
\lean{AlgebraicCombinatorics.FPS.isInvertible_iff_isUnit}
\leanhelper
\leanok
The definition of invertibility from Definition~\ref{def.commring.inverse} is equivalent
to the definition of a unit in Mathlib.
\end{lemma}

\begin{proof}
\leanok
Both express the existence of a multiplicative inverse; the equivalence is
immediate from the definitions.
\end{proof}

\begin{lemma}
\label{lem.commring.invertible-mul}
\lean{AlgebraicCombinatorics.FPS.isInvertible_mul}
\leanhelper
\leanok
If $a$ and $b$ are invertible in $L$, then $a \cdot b$ is invertible.
\end{lemma}

\begin{proof}
\leanok
Convert to the equivalent notion of a unit and use the fact that the product
of two units is a unit.
\end{proof}

\begin{lemma}
\label{lem.commring.field-isUnit}
\lean{AlgebraicCombinatorics.FPS.field_isUnit_iff}
\leanhelper
\leanok
In a field $F$, an element $a$ is invertible if and only if $a \neq 0$.
\end{lemma}

\begin{proof}
\leanok
Every nonzero element of a field has a multiplicative inverse.
\end{proof}

\subsection{Fractions and integer powers}

\begin{definition}
\label{def.commring.fraction}
\lean{AlgebraicCombinatorics.FPS.fraction}
\leanhelper
\leanok
For an invertible element $a$ and any $b \in L$,
the \emph{fraction} $b/a$ is defined as $b \cdot a^{-1}$.
\end{definition}

\begin{lemma}
\label{lem.commring.fraction-eq-iff}
\lean{AlgebraicCombinatorics.FPS.fraction_eq_iff}
\leanhelper
\leanok
For an invertible element $a$ and elements $b, c \in L$, we have $b / a = c$ if and only if
$b = c \cdot a$.
\end{lemma}

\begin{proof}
\leanok
Multiply both sides by $a$ (or $a^{-1}$) and use $a \cdot a^{-1} = 1$.
\end{proof}

\begin{lemma}
\label{lem.commring.unit-zpow-neg}
\lean{AlgebraicCombinatorics.FPS.unit_zpow_neg}
\leanhelper
\leanok
For an invertible element $a$ and a natural number $n$, we have $a^{-n} = (a^{-1})^n$.
This defines negative integer powers.
\end{lemma}

\begin{proof}
\leanok
By the definition of integer powers on units.
\end{proof}

\begin{lemma}
\label{lem.commring.unit-zpow-add}
\lean{AlgebraicCombinatorics.FPS.unit_zpow_add}
\leanhelper
\leanok
For an invertible element $a$ and integers $m, n \in \mathbb{Z}$, we have
$a^{m+n} = a^m \cdot a^n$.
\end{lemma}

\begin{proof}
\leanok
Standard property of integer powers.
\end{proof}
